\section{Trace-driven HIL Validation Plan}\label{sec:hil}

\subsection{Objectives}
The HIL plan validates the observation extraction and state estimation pipeline
in a laboratory setting without requiring live satellite passes.  It is designed
to demonstrate that Doppler-time features can be extracted from real
off-the-shelf D2S hardware and processed through the LEO-DTF estimator.

\textbf{No real satellite over-the-air validation is claimed in the current
repository state.}

\subsection{Hardware Roles}

\begin{itemize}
  \item \textbf{LR1121/STM32 packet source}: The Semtech LR1121 transceiver is
        configured in transmit-only mode (LR-FHSS or LoRa) using an STM32
        host microcontroller.  The LR1121 does \emph{not} receive; it serves
        solely as a characterizable packet waveform source with known timing.
  \item \textbf{USRP B210}: A Software-Defined Radio (SDR) front-end captures
        complex baseband IQ samples at a configurable sample rate
        (recommended 1--10~MS/s).  The USRP acts as a wideband receiver to
        capture the transmitted packets.
  \item \textbf{Offline processing host}: IQ samples are stored to disk and
        processed offline.  No real-time processing is required.
\end{itemize}

\subsection{Signal Flow and Injection}
Doppler shift, CFO, drift, and propagation delay are injected \emph{offline}
into the captured IQ samples to simulate arbitrary satellite geometries and
oscillator conditions.  This trace-driven approach enables reproducible
experiments across different pass geometries without physical satellite
visibility.

The offline injector (\texttt{src/leodtf/iq\_doppler\_injector.py}) is currently
a placeholder; the HIL plan defines the intended signal model but the
implementation is left to future work.

\subsection{Validation Scope}
The HIL validates:
\begin{itemize}
  \item Packet detection and timing extraction from IQ samples.
  \item CFO and drift estimation from frequency domain analysis.
  \item End-to-end pipeline: IQ $\rightarrow$ extracted observations
        $\rightarrow$ posterior heatmap.
\end{itemize}
The HIL does \emph{not} validate:
\begin{itemize}
  \item Real satellite Doppler extraction.
  \item Outdoor antenna and RF path effects.
  \item Real-world timing synchronization quality.
\end{itemize}

\subsection{Required Artifacts}
The following outputs shall be produced for each HIL run and stored under
\texttt{data/hil\_runs/} with a datestamp:
\begin{enumerate}
  \item \texttt{capture\_IQ\_\{run\_id\}.sigmf-data} and
        \texttt{capture\_IQ\_\{run\_id\}.sigmf-meta} --- raw IQ recording.
  \item \texttt{metadata.json} --- run parameters (Tx frequency, sample rate,
        capture duration, injected Doppler, CFO, delay, SNR).
  \item \texttt{extracted\_observations.csv} --- extracted frequency and
        timestamp observations per packet.
  \item \texttt{posterior\_result.json} --- grid posterior, MAP estimate,
        HPD region, ambiguity score.
  \item \texttt{analysis.ipynb} or \texttt{analysis.py} --- notebook or script
        reproducing the posterior from the extracted observations.
\end{enumerate}

\subsection{Minimal Experiment Checklist}
Before claiming HIL validation the following must hold:
\begin{itemize}
  \setlength\itemsep{0em}
  \item $\ge 3$ independent IQ captures with different injected geometries.
  \item MAP estimate falls within the 95\% HPD region of the true injected
        position in all runs.
  \item Posterior entropy decreases monotonically with increasing observation
        SNR in synthetic baseline runs.
  \item Analysis script reproduces the posterior heatmap from
        \texttt{extracted\_observations.csv} without manual intervention.
\end{itemize}